\subsection{An idle factory does not make its remaining stock more valuable}
\label{sec:production-capacity-cost}

Producing fewer goods leaves the factory's fixed costs spread over fewer
units, but that arithmetic does not establish their inventory cost. Consider
a factory that normally makes 100 good units each year after allowing for
planned maintenance. Its annual fixed production overhead is 180. If it
makes only 60 units, dividing all 180 by 60 would assign 3 to each unit,
rather than the normal 1.8. The unsold goods would then carry part of the
cost of idle capacity into the next year.

AASB 102.13 prevents that result. Fixed production overhead is allocated
using normal capacity; low output does not increase the amount per unit,
and unallocated overhead is expensed. Abnormally high output requires a
lower allocation per unit so that inventory is not measured above cost
\citep{aasb102compiled2021}. Normal capacity is an estimate of average
production under normal conditions, including planned maintenance. It
requires evidence separate from the actual output or desired margin.

Use a synthetic manufacturer in calendar 2024 and the retained AASB 102
compilation dated 31 December 2021. Its documented application interval
includes that year; subsequent amendment and actual-entity admission remain
separate. The company makes one ordinary manufactured product, with no
opening stock, work in progress, joint products or normal production loss.
One material kit produces one good unit. A kit costs 2; direct labor costs
1 and variable production overhead costs 1 per good unit, all paid in cash.
These goods are outside the agricultural, financial-instrument and
broker-trader exceptions. Periodic weighted average is used for this
homogeneous batch under paragraphs 25--27.

The fixed 180 consists of cash factory management and maintenance of 120
and depreciation of 60. Opening factory PPE is 600, with ten remaining
years, zero residual value and admitted uniform time-based consumption.
There is one component and recoverability has been separately established.
The depreciation method therefore remains unchanged across the output
scenarios. Paragraph 12 includes factory depreciation and production
administration among production overheads; AASB 116.48--49 allows
depreciation to enter another asset's carrying amount when appropriate
\citep{aasb116dated}. Recording it in inventory does not make it a second
cash payment.

Keep excluded costs visible. A separately documented accident destroys six
unused kits costing 12 before labor or conversion. They are additional to
the kits used for good output and have no salvage value. Unrelated
head-office administration of 10 and period selling services of 8 are also
paid in cash. The abnormal loss and these unrelated expenses do not bring
the good units to their present location and condition, so paragraphs
10--16 exclude them from inventory. Factory management and head-office
administration cannot be classified solely from the word \emph{administration}.

All completed goods have supported net realisable value of 10 per unit,
with no further completion or sale costs. The example therefore still
applies the lower-of-cost-and-NRV requirement rather than suspending it.
Delivered units sell for cash of 10 each under an approved, enforceable
single-product contract: rights and payment terms are identified,
collection is probable, the contract has commercial substance, and
delivery, acceptance, title and control transfer coincide. There are no
remaining promises, returns or concessions. These facts admit the
existing AASB 15.9 and 31--38 delivery method \citep{aasb15dated}; cash
receipt by itself would not do so.

\subsubsection{Allocate the cost that production can carry}

Let \(Q\) denote good completed units and \(S\) the units sold, with
\(0\leq S\leq Q\). For this single cost pool, normal capacity is 100
and fixed overhead is 180. Each good unit first receives its direct and
variable cost of \(2+1+1=4\). When output is low, the fixed allocation
is 1.8 per unit; when output is abnormally high, the total allocation is
capped at the 180 actually incurred. For \(Q>0\), this gives
\[
 f(Q)=\frac{180}{\max(100,Q)},\qquad
 u(Q)=4+f(Q),\qquad
 U(Q)=180-Qf(Q),
\]
where \(f\) is fixed overhead per unit, \(u\) is admitted unit cost
and \(U\) is unallocated overhead expensed now. This compact formula
implements the stated homogeneous, constant-cost example. Multi-product
pools, work in progress or changed production resources require their own
allocation. With no output, there is no unit cost to divide by: stock is
zero and all 180 is current expense.

At low output \(Q=60\), the good units carry fixed overhead of
\(60\times1.8=108\), leaving 72 as current expense. Unit cost is 5.8.
Selling 40 units removes inventory cost 232; the other 20 remain at 116.
Abnormal waste 12 and the unrelated administration/selling cost 18 are
also current expenses. Revenue 400 therefore leaves pretax profit
\(400-232-72-12-18=66\). Normal output of 100 absorbs all 180 at
the same unit cost of 5.8. Abnormally high output of 150 instead allocates
\(180/150=1.2\) per unit, giving unit cost 5.2.

The comparison below changes output and sales explicitly. All arms retain
the same fixed costs and the separately stated accident. They do not have
equal total production spending: direct and variable costs change with
good output. Sold-stock cost means the carrying amount removed on sale;
AASB 102.38 also includes unallocated overhead and abnormal production
costs in the broader inventory expense.

\begin{center}
\small
\begin{tabular}{lrrrr}
\toprule
& Low & Normal & High & Idle\\
\midrule
Good units / sold units & 60 / 40 & 100 / 80 & 150 / 120 & 0 / 0\\
Unit cost & 5.8 & 5.8 & 5.2 & n/a\\
Unallocated fixed expense & 72 & 0 & 0 & 180\\
Sold-stock cost & 232 & 464 & 624 & 0\\
Closing inventory & 116 & 116 & 156 & 0\\
Broader inventory expense & 316 & 476 & 636 & 192\\
Pretax profit & 66 & 306 & 546 & $-210$\\
Closing cash & 1010 & 1250 & 1450 & 850\\
Closing PPE & 540 & 540 & 540 & 540\\
Pretax net assets & 1666 & 1906 & 2146 & 1390\\
\bottomrule
\end{tabular}
\end{center}

Opening cash is 1000 and PPE is 600, so opening net assets are 1600.
There are no payables, borrowing or owner movements in this slice. Tax
effects are uncomputed: the last row is a pretax net-assets reconciliation,
not after-tax equity, distributable profit or a complete statutory set of
statements. This boundary avoids manufacturing tax law to finish the
example. Each arm has separately reconciled cash, stock and PPE; an
after-tax extension must supply the relevant tax bases, rates and timing.

\subsubsection{Preserve both the cost origin and the cash movement}

The low-output journal path starts with 66 kits costing 132. Debit raw
materials and credit cash 132. The accident debits abnormal-loss expense
12 and credits raw materials 12. Transferring the 60 good kits to
production moves raw-material cost 120 into finished goods; direct labor
60 and variable overhead 60 enter the same goods account against cash.
No raw materials remain, and the goods initially carry 240.

The fixed costs first enter a production clearing account: debit it 120
against cash and 60 against accumulated depreciation. Then debit
finished goods 108 and unallocated-overhead expense 72, crediting the
clearing account 180. Its balance is now zero, and finished goods cost
348. The separate administration and selling entry debits expense 18
against cash. On delivery, debit cash and credit revenue 400, then debit
sold-stock expense and credit finished goods 232. The final 116 is the
cost of 20 unsold goods, not a residual chosen to make the books balance.

Cash receipts of 400 less payments of \(132+60+60+120+18=390\)
increase cash by 10. The whole factory's depreciation is 60 even though
part remains in inventory. To reconcile pretax profit with operating cash,
add that noncash consumption and subtract the full inventory increase:
\[
 \Delta\mathrm{cash}=66+60-116=10.
\]
The stock includes depreciation of
\(20\times60/100=12\). An equivalent reconciliation adds only the
48 of depreciation reaching current profit, but then subtracts only the
cash-cost part of stock, \(116-12=104\): \(66+48-104=10\).
Adding back 48 while subtracting all 116 would wrongly report a cash
decrease of 2. The two valid methods differ in presentation, not in the
underlying cost population.

Selling all 60 good units in the low-output arm provides a second check.
Production and its payments are unchanged. Revenue becomes 600,
sold-stock expense 348, profit 150 and cash 1210; inventory becomes zero.
The additional 200 cash receipt releases inventory cost 116 and increases
profit by 84. Total pretax net assets become 1750. This is why a forecast
needs both production and sell-through, rather than an inventory margin
inferred from total factory spending.

\subsubsection{Balanced books can still defer the wrong costs}

If low output absorbed all fixed overhead using \(180/60=3\) per
unit, unit cost would be 7. Unsold stock would be 140 and profit 90,
each 24 above the admitted result; cash would remain 1010. The books
would balance, but they would postpone idle-capacity expense in the
remaining goods. At high output, applying 1.8 without a cap would allocate
270 of fixed cost when only 180 was incurred. Crediting a variance of 90
to balance that over-allocation would leave stock 174 and profit 564,
both overstated by 18. A balanced variance cannot create an incurred cost.

Capitalizing abnormal loss 12 and unrelated period costs 18 into the
60-unit batch would defer \(30\times20/60=10\) of excluded expense
in unsold stock. This is a separate classification error from the
capacity error. Each resource therefore retains a purpose and identity,
its membership in the cost pool, and its cash or noncash origin. The
state also retains capacity evidence, good and lost units, batch cost
and quantity, and the specific sale removals. A duplicate resource or
negative stock is invalid; a revised capacity estimate creates a new
supported allocation decision instead of overwriting an earlier result.

Unknown normal capacity leaves the fixed allocation unresolved. Unknown
NRV prevents a complete closing measurement, and missing control-transfer
facts prevent revenue recognition. Work in progress, joint products and
heterogeneous cost populations require a different specified model rather
than silent use of this unit formula. The proposed procedure and checked
case define those boundaries; they do not establish production-engine
conformance or an actual manufacturer's accounting treatment.
